\paragraph{Summary.}
This work introduces \ouralgo, an evolutionary framework addressing critical limitations in LLM-driven scientific discovery through improved sample efficiency and open-source accessibility.
\ouralgo achieves state-of-the-art results across four domains: circle packing with 150 evaluations (orders of magnitude improvement), sophisticated AIME reasoning scaffolds, ALE-Bench algorithmic improvements, and novel mixture-of-expert load balancing.
%
\paragraph{Limitations.}
Our implementation uses fixed configurations with limited automatic control over exploration-exploitation balance, which may vary across domains.
Task specification requires manual human expertise for objective functions and evaluation.
The framework is constrained to problems with well-defined numerical objectives, limiting its applicability to diverse evaluation domains.
%
\paragraph{Future Directions.}
Automated task specification through LLM task generation could enable greater autonomy and unlock applications in unexplored domains.
Transitioning to true open-endedness, where systems generate their own objectives, represents a compelling frontier.
Self-referential refinement and online meta-learning offer opportunities for continuously improving discovery.
%
\paragraph{Broader Impact \& Ethical Considerations.}
\ouralgo's open-source release further democratizes advanced evolutionary optimization, making it accessible to researchers and practitioners previously lacking access to proprietary systems.
The framework's exceptional sample efficiency reduces computational barriers for resource-constrained environments.
However, API costs from large-scale LLM usage could create economic barriers, potentially constraining democratization goals.